\section{The \texttt{Undulator} McXtrace Component}
Model of an undulator source

\subsection*{Identification}
\begin{itemize}
  \item \textbf{Author:} Erik B. Knudsen
  \item \textbf{Origin:} DTU Physics
  \item \textbf{Date:} May, 2013.
\end{itemize}

\subsection*{Description}
A undulator source model based on the derivation by K.J. Kim, AIP, conf. proc., 184, 1989. doi:10.1063/1.38046.

SOLEIL\_PX2a U24

\begin{verbatim}
Example: Undulator( E0=12.65, dE=1, Ee=2.75, dEe=0.001, Ie=0.5, K=1.788, Nper=80,
\end{verbatim}

lu=24e-3, sigey=9.3e-6, sigex=215.7e-6, sigepx=29.3e-6, sigepy=4.2e-6,

\begin{verbatim}
dist=29.5, E1st=12.400 )
\end{verbatim}

\subsection*{Input parameters}
Parameters in \textbf{boldface} are required; the others are optional.

\begin{longtable}{p{0.22\textwidth}p{0.12\textwidth}p{0.46\textwidth}p{0.14\textwidth}}
\toprule
\textbf{Name} & \textbf{Unit} & \textbf{Description} & \textbf{Default} \\
\midrule
\endhead
E0 & keV & Center of emitted energy spectrum. & 0 \\
dE & keV & Half-width of emitted energy spectrum. & 0 \\
phase & rad & Initial phase of radiation. & 0 \\
randomphase & 0/1 & If !=0 phase will be random (I.e. the emitted radiation is completely incoherent). & 1 \\
Ee & GeV & Storage ring electron energy [typically a few GeV). & 2.4 \\
dEe & percent & Relative electron energy beam spread (sigma/Ee). & 0 \\
Ie & A & Ring current. & 0.4 \\
B & T & Peak magnet field strength. Overrides K. & 0 \\
K & 1 & Dimensionless deflection undulator parameter. When K \textgreater{}\textgreater{} 1 (ie B*lu is large) you get a wiggler. & 0 \\
Nper & int & Number of magnetic periods in the undulator. & 1 \\
lu & m & Magnetic period length of the undulator aka lambda\_u. & 16e-3 \\
sigey & m & Electron ring beam size in vertical plane (rms). & 0 \\
sigex & m & Electron ring beam size in horizontal plane (rms). & 0 \\
sigepx & rad & Electron ring beam horizontal divergence (rms). & 0 \\
sigepy & rad & Electron ring beam vertical divergence (rms). & 0 \\
focus\_xw & m & Width of target window. & 0 \\
focus\_yh & m & Height of target window. & 0 \\
dist & m & Distance from source plane to target window along the optical axis. & 1 \\
quick\_integ & 0/1 & If nonzero, use faster (but less accurate) integration scheme. & 0 \\
E1st & keV & Energy of the fundmental (1st) undulator harmonic. & 0 \\
verbose & 0/1 & If nonzero, output extra information. & 0 \\
Br & T & Remanent field (1.35T for Nd2Fe14B) for gap estimate & 1.35 \\
\bottomrule
\end{longtable}

\subsection*{Links}
\begin{itemize}
  \item Component source code found in file \texttt{Undulator.comp}.
\end{itemize}
\IfFileExists{sources/Undulator_static.tex}{\input{sources/Undulator_static.tex}}{}